% =====================================================================
%  capitulos/01_intro.tex  --  SOLO CONTENIDO
% =====================================================================
%  Este fichero contiene únicamente texto. Fíjate en lo que NO hay:
%    - \documentclass, \usepackage ni \begin{document}: eso es de main.tex
%      y del .sty. main.tex lo incorpora con % =====================================================================
%  capitulos/01_intro.tex  --  SOLO CONTENIDO
% =====================================================================
%  Este fichero contiene únicamente texto. Fíjate en lo que NO hay:
%    - \documentclass, \usepackage ni \begin{document}: eso es de main.tex
%      y del .sty. main.tex lo incorpora con % =====================================================================
%  capitulos/01_intro.tex  --  SOLO CONTENIDO
% =====================================================================
%  Este fichero contiene únicamente texto. Fíjate en lo que NO hay:
%    - \documentclass, \usepackage ni \begin{document}: eso es de main.tex
%      y del .sty. main.tex lo incorpora con % =====================================================================
%  capitulos/01_intro.tex  --  SOLO CONTENIDO
% =====================================================================
%  Este fichero contiene únicamente texto. Fíjate en lo que NO hay:
%    - \documentclass, \usepackage ni \begin{document}: eso es de main.tex
%      y del .sty. main.tex lo incorpora con \include{capitulos/01_intro}.
%    - colores ni márgenes: eso es diseño y vive en styles/mistilo.sty.
%
%  PISTAS PARA EDITAR
%    - Una línea en blanco separa párrafos.
%    - ~ es un espacio que NO permite salto de línea ("sección~2").
%    - \emph{...} = cursiva de énfasis;  \textbf{...} = negrita.
% =====================================================================

% \section crea una sección numerada y la añade al índice.
\section{Introducción}

La programación paralela es hoy una herramienta cotidiana, pero escribir
código paralelo eficiente y portable sigue siendo difícil. Una línea de
trabajo de la Universidad de Valladolid propone describir los programas
con especificaciones de alto nivel y dejar que un compilador modular
genere el código final~\cite{gonzalez2011trasgo}.

% \cite{clave} inserta una cita. La clave "gonzalez2011trasgo" apunta a
% una entrada de bib/refs.bib. Si la clave no existe o está mal escrita
% saldrá [?] y un aviso al compilar. El formato visible ([1],
% (González, 2011)...) lo decide el estilo elegido en main.tex, no este
% fichero.
%
% Para citar VARIAS referencias a la vez:
%     \cite{gonzalez2011trasgo,fernandez2013hdt}

% Entorno "aviso" definido en mistilo.sty. Sin argumento, el título
% por defecto es "Aviso".
\begin{aviso}
  Este documento es solo un ejemplo para el taller: las referencias son
  reales, el resto es ficticio.
\end{aviso}

% \subsection crea una subsección (1.1, 1.2...).
\subsection{Estructura del informe}
% \ref{etiqueta} pone el número de algo marcado con \label{etiqueta}.
% Aquí "sec:datos" está definida en capitulos/02_datos.tex. Si se
% reordenan las secciones, el número se actualiza solo (tras 2 pasadas).
La sección~\ref{sec:datos} describe un formato binario para publicar
grandes conjuntos de datos enlazados~\cite{fernandez2013hdt} y
presenta una tabla resumen.
.
%    - colores ni márgenes: eso es diseño y vive en styles/mistilo.sty.
%
%  PISTAS PARA EDITAR
%    - Una línea en blanco separa párrafos.
%    - ~ es un espacio que NO permite salto de línea ("sección~2").
%    - \emph{...} = cursiva de énfasis;  \textbf{...} = negrita.
% =====================================================================

% \section crea una sección numerada y la añade al índice.
\section{Introducción}

La programación paralela es hoy una herramienta cotidiana, pero escribir
código paralelo eficiente y portable sigue siendo difícil. Una línea de
trabajo de la Universidad de Valladolid propone describir los programas
con especificaciones de alto nivel y dejar que un compilador modular
genere el código final~\cite{gonzalez2011trasgo}.

% \cite{clave} inserta una cita. La clave "gonzalez2011trasgo" apunta a
% una entrada de bib/refs.bib. Si la clave no existe o está mal escrita
% saldrá [?] y un aviso al compilar. El formato visible ([1],
% (González, 2011)...) lo decide el estilo elegido en main.tex, no este
% fichero.
%
% Para citar VARIAS referencias a la vez:
%     \cite{gonzalez2011trasgo,fernandez2013hdt}

% Entorno "aviso" definido en mistilo.sty. Sin argumento, el título
% por defecto es "Aviso".
\begin{aviso}
  Este documento es solo un ejemplo para el taller: las referencias son
  reales, el resto es ficticio.
\end{aviso}

% \subsection crea una subsección (1.1, 1.2...).
\subsection{Estructura del informe}
% \ref{etiqueta} pone el número de algo marcado con \label{etiqueta}.
% Aquí "sec:datos" está definida en capitulos/02_datos.tex. Si se
% reordenan las secciones, el número se actualiza solo (tras 2 pasadas).
La sección~\ref{sec:datos} describe un formato binario para publicar
grandes conjuntos de datos enlazados~\cite{fernandez2013hdt} y
presenta una tabla resumen.
.
%    - colores ni márgenes: eso es diseño y vive en styles/mistilo.sty.
%
%  PISTAS PARA EDITAR
%    - Una línea en blanco separa párrafos.
%    - ~ es un espacio que NO permite salto de línea ("sección~2").
%    - \emph{...} = cursiva de énfasis;  \textbf{...} = negrita.
% =====================================================================

% \section crea una sección numerada y la añade al índice.
\section{Introducción}

La programación paralela es hoy una herramienta cotidiana, pero escribir
código paralelo eficiente y portable sigue siendo difícil. Una línea de
trabajo de la Universidad de Valladolid propone describir los programas
con especificaciones de alto nivel y dejar que un compilador modular
genere el código final~\cite{gonzalez2011trasgo}.

% \cite{clave} inserta una cita. La clave "gonzalez2011trasgo" apunta a
% una entrada de bib/refs.bib. Si la clave no existe o está mal escrita
% saldrá [?] y un aviso al compilar. El formato visible ([1],
% (González, 2011)...) lo decide el estilo elegido en main.tex, no este
% fichero.
%
% Para citar VARIAS referencias a la vez:
%     \cite{gonzalez2011trasgo,fernandez2013hdt}

% Entorno "aviso" definido en mistilo.sty. Sin argumento, el título
% por defecto es "Aviso".
\begin{aviso}
  Este documento es solo un ejemplo para el taller: las referencias son
  reales, el resto es ficticio.
\end{aviso}

% \subsection crea una subsección (1.1, 1.2...).
\subsection{Estructura del informe}
% \ref{etiqueta} pone el número de algo marcado con \label{etiqueta}.
% Aquí "sec:datos" está definida en capitulos/02_datos.tex. Si se
% reordenan las secciones, el número se actualiza solo (tras 2 pasadas).
La sección~\ref{sec:datos} describe un formato binario para publicar
grandes conjuntos de datos enlazados~\cite{fernandez2013hdt} y
presenta una tabla resumen.
.
%    - colores ni márgenes: eso es diseño y vive en styles/mistilo.sty.
%
%  PISTAS PARA EDITAR
%    - Una línea en blanco separa párrafos.
%    - ~ es un espacio que NO permite salto de línea ("sección~2").
%    - \emph{...} = cursiva de énfasis;  \textbf{...} = negrita.
% =====================================================================

% \section crea una sección numerada y la añade al índice.
\section{Introducción}

La programación paralela es hoy una herramienta cotidiana, pero escribir
código paralelo eficiente y portable sigue siendo difícil. Una línea de
trabajo de la Universidad de Valladolid propone describir los programas
con especificaciones de alto nivel y dejar que un compilador modular
genere el código final~\cite{gonzalez2011trasgo}.

% \cite{clave} inserta una cita. La clave "gonzalez2011trasgo" apunta a
% una entrada de bib/refs.bib. Si la clave no existe o está mal escrita
% saldrá [?] y un aviso al compilar. El formato visible ([1],
% (González, 2011)...) lo decide el estilo elegido en main.tex, no este
% fichero.
%
% Para citar VARIAS referencias a la vez:
%     \cite{gonzalez2011trasgo,fernandez2013hdt}

% Entorno "aviso" definido en mistilo.sty. Sin argumento, el título
% por defecto es "Aviso".
\begin{aviso}
  Este documento es solo un ejemplo para el taller: las referencias son
  reales, el resto es ficticio.
\end{aviso}

% \subsection crea una subsección (1.1, 1.2...).
\subsection{Estructura del informe}
% \ref{etiqueta} pone el número de algo marcado con \label{etiqueta}.
% Aquí "sec:datos" está definida en capitulos/02_datos.tex. Si se
% reordenan las secciones, el número se actualiza solo (tras 2 pasadas).
La sección~\ref{sec:datos} describe un formato binario para publicar
grandes conjuntos de datos enlazados~\cite{fernandez2013hdt} y
presenta una tabla resumen.
